\section{Local Gaussian analysis}\label{sec:gaussian}

Recall from \eqref{3.3} that
\[
L P^{\mathrm{quot}}=\sum_{h\bmod L}F(h).
\tag{6.1}
\]
We now evaluate this coefficient on its natural scale $1/\sigma$.

Fix $C\ge1$. The limit $X\to\infty$ will be taken with $C$ fixed, and only afterwards will $C\to\infty$. Put
\[
N=\left\lfloor\frac{C}{\sigma}\right\rfloor.
\tag{6.2}
\]
Since $\sigma^{-1}\asymp_bX\log X$, for fixed $C$ and large $X$,
\[
N<X^2/4<M_{\mathrm{dec}}.
\tag{6.3}
\]

By Section \ref{sec:anchor}, every top assignment is either noncoherent, with anchor energy at least $F_{\mathcal B}$, or coherent with a unique exact integer label $m$. For a coherent assignment, exactly one of the ranges
\[
\begin{array}{ll}
\mathrm{I}.& |m|\le N,\\
\mathrm{II}.& N<|m|\le X^2/4,\\
\mathrm{III}.& X^2/4<|m|\le M_{\mathrm{dec}},\\
\mathrm{IV}.& |m|>M_{\mathrm{dec}}
\end{array}
\tag{6.4}
\]
holds. Lemma \ref{lem:weighted-skeleton} then separates, within each anchor fibre, the decoded skeleton from all other row states. This partitions the full Fourier sum into a Gaussian core, three coherent tails, and the noncoherent/nondecoder remainder.

\subsection{The Gaussian core}

For $|m|\le N$, Propositions \ref{prop:lower-decoder} and \ref{prop:target-decoder} identify the decoded skeleton with the genuine integer frequency $h=m\pmod L$. Since $\theta$ stays in a compact subinterval of $(0,1)$, there are constants $\rho_b>0$ and $M_b<\infty$ such that, uniformly for $|z|\le\rho_b$,
\[
\log\bigl((1-\theta)+\theta e(z)\bigr)
=
2\pi i\theta z
-2\pi^2\theta(1-\theta)z^2
+R_\theta(z),
\tag{6.5}
\]
with
\[
|R_\theta(z)|\le M_b|z|^3.
\tag{6.6}
\]
Every denominator in $\mathcal E$ is at least $X^2$ for large $X$. Uniformly for $|m|\le C/\sigma$,
\[
\max_{e\in\mathcal E}\frac{|m|}{e}
\ll_C\frac{\log X}{X}
\longrightarrow0.
\tag{6.7}
\]
Moreover
\[
\sum_{e\in\mathcal E}\left|\frac{m}{e}\right|^3
\le
\frac{|m|^3}{X^2}
\sum_{e\in\mathcal E}\frac1{e^2}
\ll_b
\frac{C^3}{\sigma X^2}
\ll_b
C^3\frac{\log X}{X}
\longrightarrow0.
\tag{6.8}
\]
Thus the aggregate cubic remainder in \eqref{6.5} is $o_C(1)$ uniformly throughout the major range.

The linear term cancels exactly by \eqref{2.11}:
\[
2\pi i m\theta\sum_{e\in\mathcal E}\frac1e
-\frac{2\pi i m}{b}=0.
\tag{6.9}
\]
The quadratic term is $-2\pi^2m^2\sigma^2$. Therefore
\[
F(m)
=
\exp\bigl(-2\pi^2m^2\sigma^2+\varepsilon_{X,m}\bigr),
\qquad
\max_{|m|\le N}|\varepsilon_{X,m}|
\longrightarrow0
\tag{6.10}
\]
for every fixed $C$.

\begin{proposition}[Normalized major arc]\label{prop:normalized-major}
For every fixed $C>0$,
\[
\sigma
\sum_{|m|\le C/\sigma}F(m)
\longrightarrow
\int_{-C}^{C}e^{-2\pi^2u^2}\,du.
\tag{6.11}
\]
\end{proposition}

\begin{proof}
Put $u_m=m\sigma$. The mesh is $\sigma\to0$, and \eqref{6.10} is uniform on $[-C,C]$. Thus the left side is a Riemann sum for the integral in \eqref{6.11}.
\end{proof}

The Gaussian core is therefore of order $1/\sigma$. We compare every remaining contribution with this same scale.

\subsection{The first coherent tail}

For
\[
N<|m|\le X^2/4,
\tag{6.12}
\]
the decoded skeleton is still the genuine integer frequency $m$. Since every denominator is at least $X^2$,
\[
|m|/e\le1/4,
\]
so $\|m/e\|=|m|/e$. Lemma \ref{lem:bernoulli-modulus} gives
\[
|F(m)|\le e^{-c_bm^2\sigma^2}.
\tag{6.13}
\]
A Gaussian integral comparison yields
\[
\sum_{|m|>C/\sigma}e^{-c_bm^2\sigma^2}
\le
K_b\sigma^{-1}e^{-c_bC^2/2}.
\tag{6.14}
\]
Hence
\[
\limsup_{X\to\infty}
\sigma
\sum_{N<|m|\le X^2/4}|F(m)|
\le
K_be^{-c_bC^2/2}.
\tag{6.15}
\]
This is the only complementary range whose normalized bound need not tend to zero at fixed $C$; it will disappear when $C\to\infty$.

\subsection{The intermediate coherent tail}

Suppose
\[
X^2/4<|m|\le M_{\mathrm{dec}}.
\tag{6.16}
\]
All prime coordinates still decode to $m$. By \eqref{5.25},
\[
|F(h_m)|
\le
\exp\!\left(-c\frac{|m|}{(\log|m|)^2}\right).
\tag{6.17}
\]
Throughout \eqref{6.16}, $\log|m|\asymp\log X$. There are at most $2M_{\mathrm{dec}}+1=O(X^4)$ labels, so
\[
\sum_{X^2/4<|m|\le M_{\mathrm{dec}}}|F(h_m)|
\ll
X^4\exp\!\left(-c\frac{X^2}{\log^2X}\right)
=o(\sigma^{-1}).
\tag{6.18}
\]

\subsection{Large coherent anchor labels}

For a coherent top label $m$, the internal anchor factor is at most
\[
\exp(-\kappa_bm^2\sigma_{\mathcal B,0}^2).
\]
All row and residual factors have modulus at most one. Thus
\[
\sum_{|m|>M_{\mathrm{dec}}}|F(h_m)|
\ll
\sigma_{\mathcal B,0}^{-1}
\exp\!\left(-cM_{\mathrm{dec}}^2\sigma_{\mathcal B,0}^2\right).
\tag{6.19}
\]
From \eqref{4.38} and \eqref{5.17},
\[
M_{\mathrm{dec}}^2\sigma_{\mathcal B,0}^2
\gg
\frac{X^2}{\log^6Z},
\tag{6.20}
\]
while $\sigma_{\mathcal B,0}^{-1}=O(Z\log Z)$. Hence
\[
\sum_{|m|>M_{\mathrm{dec}}}|F(h_m)|
=o(\sigma^{-1}).
\tag{6.21}
\]

\subsection{Incoherent anchors and nondecoder states}

By \eqref{4.44}, the total anchor modulus of noncoherent assignments is
\[
O(e^{-cF_{\mathcal B}}).
\tag{6.22}
\]
By \eqref{5.14}, the complete contribution of nondecoder row states, over coherent and noncoherent anchors alike, is
\[
O\!\left(
Z\log Z\cdot
Z^2e^{-cZ/\log^3Z}
\right).
\tag{6.23}
\]
Both \eqref{6.22} and \eqref{6.23} are $o(\sigma^{-1})$, because $\sigma^{-1}\asymp_bX\log X$ and $Z=X^3$.

\subsection{Assembly and the local limit}

Let $\mathfrak R_{X,C}$ be the full Fourier contribution outside the genuine integer range $|m|\le C/\sigma$. Equations \eqref{6.15}, \eqref{6.18}, \eqref{6.21}, \eqref{6.22}, and \eqref{6.23} give
\[
\limsup_{X\to\infty}
\sigma|\mathfrak R_{X,C}|
\le
K_be^{-c_bC^2/2}.
\tag{6.24}
\]
Combining this with Proposition \ref{prop:normalized-major},
\[
\limsup_{X\to\infty}
\left|
\sigma L P^{\mathrm{quot}}
-
\int_{-C}^{C}e^{-2\pi^2u^2}\,du
\right|
\le
K_be^{-c_bC^2/2}.
\tag{6.25}
\]
Now let $C\to\infty$. Since
\[
\int_{\R}e^{-2\pi^2u^2}\,du
=
\frac1{\sqrt{2\pi}},
\tag{6.26}
\]
we obtain
\[
\sigma L P^{\mathrm{quot}}
\longrightarrow
\frac1{\sqrt{2\pi}}.
\tag{6.27}
\]
By Lemma \ref{lem:no-wrap}, $P^{\mathrm{quot}}=P$. Therefore
\[
L P
\sim
\frac1{\sqrt{2\pi}\,\sigma}.
\tag{6.28}
\]
This is the central value \eqref{1.7} of Theorem \ref{thm:local-profile}.

The order of limits is part of the argument:
\[
X\to\infty\quad\text{at fixed }C,
\qquad\text{then}\qquad
C\to\infty.
\tag{6.29}
\]
No diagonal choice $C=C(X)\to\infty$ is asserted.

\begin{corollary}[Quantitative positivity]\label{cor:positive}
There is $c_b>0$ such that, for all sufficiently large $X$,
\[
LP\ge\frac{c_b}{\sigma}.
\tag{6.30}
\]
In particular $P>0$.
\end{corollary}

The stronger asymptotic \eqref{6.28}, rather than mere positivity, will drive the entropy-scale conclusions of Section \ref{sec:multiplicity}. The full moving-target form of Theorem \ref{thm:local-profile} is proved in Section \ref{sec:moving-target} by inserting one unit character into the same analysis.
